\section{Conclusion}
In synchronous annealing, the previous state of each spin returns to it through its neighbors, an echo that earlier parallel annealers suppress with tuned pinning or time-step couplings. Anechoic subtracts the echo: dynamical mean-field theory gives its coefficient, which the engine counts with one population count per step (Onsager-online) or approximates from the temperature (Onsager-$\kappa T$). At every step budget on K2000, both forms reach lower Ising energy than five published binary-spin algorithms run at their papers' settings. On the AMD V80, twelve flip-proportional engines (\cycStep{} cycles per step plus \cycFlip{} per flip) turn the steps and flips that the correction saves into time and reach the K2000 target in \ttsPr{}: \spdBSB{} faster than the fastest published hardware~\cite{goto2021high} and $4.6\times$ faster than STATICA's published schedules on the same board. With 2-bit couplings, Onsager-$\kappa T$ is also faster than plain SCA on 48 of 51 G-set instances. The same bit-exact engine needs \gpuSpd{} more time and \gpuEratio{} more energy per solution on an NVIDIA GH200, though the GPU's 240 concurrent anneals give 2.3$\times$ higher throughput. As a GlobalFoundries 22\,nm ASIC at 1.25\,GHz, twelve engines would reach the target in about 10\,\textmu s, $26\times$ faster than the fastest published hardware, with 0.20\,mJ of engine energy per solution.

\section*{Use of AI Tools}\label{subsec:ai}
Artificial intelligence (AI) tools such as Gemini~3.1~Pro, Claude Opus~5.5 and GPT-6 Astra are used to proofread the manuscript and to check the related work and the theory, including the machine-checked proofs of Section~\ref{subsec:echo}. Additionally, they are used to re-implement the published algorithms used for comparison (Section~\ref{subsec:algcomp}), while all values reported by other papers, including those in Table~\ref{tab:comp}, are quoted directly from the cited publications.
